n este trabajo se aborda el problema AniMarathon, cuyo objetivo consiste en seleccionar capítulos de distintos animes maximizando la satisfacción total del usuario, respetando restricciones de tiempo y energía. El problema presenta la particularidad de que sólo es posible visualizar prefijos de cada anime, pudiendo obtener bonificaciones adicionales al completar una serie.

Para resolver el problema se implementaron cuatro enfoques: fuerza bruta, dos algoritmos greedy y programación dinámica. Además, se realizaron experimentos computacionales utilizando instancias generadas automáticamente de distintos tamaños con el fin de analizar la calidad de las soluciones y el rendimiento computacional de cada método.

Los resultados muestran que la fuerza bruta resulta inviable para instancias grandes debido a su complejidad exponencial, mientras que los algoritmos greedy ofrecen tiempos de ejecución reducidos a costa de perder optimalidad. Por su parte, la programación dinámica logra obtener soluciones óptimas manteniendo tiempos de ejecución razonables, constituyéndose como la mejor alternativa práctica para resolver el problema.